\documentclass[11pt,a4paper]{article}

\usepackage[margin=1.7cm]{geometry}
\usepackage{fontspec}
\usepackage{longtable}
\usepackage{array}
\usepackage{tikz}
\usetikzlibrary{arrows.meta}

\setmainfont{Noto Serif CJK TC}
\setsansfont{Noto Sans CJK TC}
\setmonofont{DejaVu Sans Mono}
\setlength{\parindent}{0pt}
\setlength{\LTleft}{0pt}
\setlength{\LTright}{0pt}

\newcommand{\dpath}[1]{{\footnotesize\texttt{\detokenize{#1}}}}
% This TeX installation has no CJK line-breaking package. Mark natural
% break points in the long table descriptions with this command.
\newcommand{\cjkbreak}{\allowbreak}

\begin{document}

\section*{資料流程框架}

\begin{center}
\begin{tikzpicture}[
  >=Latex,
  flow/.style={->, thick},
  stage/.style={draw, rounded corners, align=center, text width=7.3cm,
    minimum height=0.82cm, inner sep=4pt, font=\small},
  side/.style={draw, rounded corners, align=center, text width=4.15cm,
    minimum height=0.82cm, inner sep=4pt, font=\small}
]
  \node[stage] (proof) at (0,0) {執行 proof};
  \node[stage] (raw) at (0,-1.3)
    {\dpath{data/raw/}\\Collector 收集並限制保留的 records};
  \node[stage] (merged) at (0,-2.6)
    {\dpath{data/merged/}\\跨 proofs 合併同一個 tactic 的 records};
  \node[stage] (dedup) at (0,-3.9)
    {\dpath{data/dedup/}\\完全去重};
  \node[stage] (sampling) at (0,-5.2)
    {\dpath{data/sampling/}\\similarity check、各組配額與空缺遞補};
  \node[stage] (generation) at (0,-6.5)
    {\dpath{data/generation/}\\準備實際要送給 LLM 的 records 和 prompt};
  \node[stage] (descriptions) at (0,-7.8)
    {\dpath{llm_output/descriptions/}\\LLM 生成的 descriptions};
  \node[stage] (judge) at (0,-9.1)
    {\dpath{metrics/judge_responses/}\\LLM judge 的回答};
  \node[stage] (record) at (0,-10.4)
    {\dpath{metrics/record_scores/}\\將回答整理成逐筆分數};
  \node[stage] (tactic) at (0,-11.7)
    {\dpath{metrics/tactic_scores/}\\按 tactic／description 彙整};

  \node[side] (selection) at (6.4,-5.2)
    {\dpath{metrics/selection/}\\篩選決定與統計};
  \node[side] (evaluation) at (6.4,-7.8)
    {\dpath{data/evaluation/}\\選定用來評分的 records};

  \draw[flow] (proof) -- (raw);
  \draw[flow] (raw) -- (merged);
  \draw[flow] (merged) -- (dedup);
  \draw[flow] (dedup) -- (sampling);
  \draw[flow] (sampling) -- (generation);
  \draw[flow] (generation) -- (descriptions);
  \draw[flow] (descriptions) -- (judge);
  \draw[flow] (judge) -- (record);
  \draw[flow] (record) -- (tactic);
  \draw[flow] (sampling) -- (selection);
  \draw[flow] (evaluation.south) |- (judge.east);
\end{tikzpicture}
\end{center}

\small
\textbf{跨階段追溯：}\dpath{data/metadata/} 記錄上述各階段的來源、設定與檔案關係。

\clearpage
\section*{各階段資料夾說明}

\setlength{\tabcolsep}{4pt}
\renewcommand{\arraystretch}{1.35}
\begin{longtable}{|>{\raggedright\arraybackslash}p{1.7cm}|>{\raggedright\arraybackslash}p{4.6cm}|>{\raggedright\arraybackslash}p{10.1cm}|}
\hline
\textbf{順序} & \textbf{資料夾} & \textbf{裡面存什麼} \\ \hline
\endfirsthead
\hline
\textbf{順序} & \textbf{資料夾} & \textbf{裡面存什麼} \\ \hline
\endhead
1 & \dpath{data/raw/} & Collector 跑完 proof 後輸出的 records。\cjkbreak 已經受到 collector 的大小、容量與選取規則限制，\cjkbreak 因此不是執行過程中的所有 records。 \\ \hline
2 & \dpath{data/merged/} & 把不同 proofs 中同一個 tactic 的候選 records 集合在一起，\cjkbreak 此階段仍可包含重複資料。 \\ \hline
3 & \dpath{data/dedup/} & 依既有 record identity 規則，\cjkbreak 移除完全重複後的 records。\cjkbreak 尚未做後面的 similarity 與最終配額選擇。 \\ \hline
4 & \dpath{data/sampling/} & 經過 similarity check、分組配額與遞補後，\cjkbreak 最後選出的 records，\cjkbreak 供 description generation 使用。 \\ \hline
4 的紀錄 & \dpath{metrics/selection/} & 為什麼保留或排除某筆 record、相似度比較、\cjkbreak 各組數量與篩選設定。\cjkbreak 這不是 LLM 對 description 的評分。 \\ \hline
5 & \dpath{data/generation/} & 實際生成請求：\cjkbreak 包含送進 LLM 的 records、prompt 和相關設定。\cjkbreak 若生成時另外抽樣或截斷，\cjkbreak 這裡記錄實際使用的版本。 \\ \hline
6 & \dpath{llm_output/descriptions/} & LLM 產生的 descriptions，\cjkbreak 以及生成嘗試與狀態等紀錄。 \\ \hline
7 & \dpath{data/evaluation/} & Judge 使用的 records、\cjkbreak 候選 descriptions 的對應資訊與評分任務。\cjkbreak 放在這裡不代表它們與生成資料沒有重疊，\cjkbreak 是否獨立要看實驗設計。 \\ \hline
8 & \dpath{metrics/judge_responses/} & 保存 judge 的回答，\cjkbreak 包括逐筆判斷、理由、用量和呼叫錯誤。 \\ \hline
9 & \dpath{metrics/record_scores/} & 把有效回答整理成每筆 record 對每份 description 的分數與理由。\cjkbreak 目前這個 judge 的對應是 full=1、partial=0.5、none=0。 \\ \hline
10 & \dpath{metrics/tactic_scores/} & 將逐筆分數彙整成每個 tactic／description 的結果，\cjkbreak 供比較、報表與表格使用。 \\ \hline
全流程 & \dpath{data/metadata/} & s2n-bignum commit、實驗設定、來源檔案、雜湊、搬移對照等追溯資訊。\cjkbreak 各資料清單的 index 也可以與 records 放在同一階段資料夾。 \\ \hline
\end{longtable}

\end{document}
